\documentclass[14pt]{extreport}
\usepackage{indentfirst}
\usepackage[T2A]{fontenc}
\usepackage[utf8]{inputenc}
\usepackage[english,russian]{babel}
\usepackage{parskip}
\usepackage[top = 20mm, left = 20mm, right = 20mm, bottom = 20mm]{geometry}
\setlength\parindent{12.5mm} 
\renewcommand{\baselinestretch}{1.5} 
\def\hrf#1{\hbox to#1{\hrulefill}} 
\usepackage{ragged2e} 
\usepackage{caption}
\usepackage{tikz}
\usepackage{circuitikz}
\usetikzlibrary{decorations.pathmorphing,arrows.meta}
\usepackage{multicol}
\usepackage{float}
\usepackage{tipa}
\usepackage{graphicx}
\usepackage{enumitem}
\usepackage{amsmath}
\usepackage{amssymb}
\usepackage{listings}
\usepackage{moreverb}
\usepackage{minted}
\setcounter{tocdepth}{3}
\setcounter{secnumdepth}{3}

\lstset{
    language=C++,
    basicstyle=\ttfamily\footnotesize,
    numbers=left,               % Нумерация строк слева
    numberstyle=\tiny\color{gray}, % Стиль номеров строк
    stepnumber=1,               % Номер на каждой строке
    numbersep=5pt,              % Расстояние от номеров до кода
    breaklines=true,            % Перенос длинных строк
    tabsize=2,                  % Размер табуляции
    extendedchars=false,
    keepspaces=true,
    showstringspaces=false,
    frame=single,               % Рамка вокруг листинга
    captionpos=t,               % Подпись сверху
    xleftmargin=10pt,
    xrightmargin=5pt
}
\usepackage[colorlinks=true, citecolor=blue, urlcolor=blue, linkcolor=blue]{hyperref}
\renewcommand\thesection{\arabic{section}}


\begin{document}

\thispagestyle{empty}
    \centering{
    \hspace{0.2cm}МИНИСТЕРСТВО НАУКИ И ВЫСШЕГО ОБРАЗОВАНИЯ РОССИЙСКОЙ ФЕДЕРАЦИИ\\
    Федеральное государственное автономное образовательное учреждение высшего образования\\
    {\bf «Санкт-Петербургский политехнический университет Петра Великого»}\\[10pt]
    Институт компьютерных наук и кибербезопасности\\
    Высшая школа технологий искусственного интелекта\\
    02.03.01 Математика и компьютерные науки\\
    \vspace{1.5cm}
    {\large
    \bf Отчет по курсовой работе\\
    \bf Структуры данных}\\
    {\large\bf Разработка игрового бота для расширенной версии игры “Крестики-нолики”}\\
    \vspace{\fill}
    }
    \vspace{1.5cm}
    \flushleft{
        {\bf Выполнил:}\\
        студент группы 5130201/50301 \hspace{2.25cm} 
        \line(1,0){120} \quad Ташлыкова Е.В.\\
        {\bf Принял:}\\
        Ассистент \hspace{66mm}
        \line(1,0){120} \quad Кирпиченко С.Р.\\
        \vspace{1cm}
        \hspace{11.6cm}<<\hrf{2em}>> \hrf{6em} 2026г.
        }
        
    \vspace{1.5cm}
    \centering{
    Санкт-Петербург, 2026
    }

\newpage
\justifying
\newpage
\renewcommand\contentsname{\vspace{-3.2cm}\centering\LargeСодержание}
\tableofcontents
\thispagestyle{empty}
\newpage
{\centering\section*{Введение}}\addcontentsline{toc}{section}{Введение}
В раках Летней Научно-исследовательской работы были использованы такие интсрумены и технологии, как VirtualBox -- это кроссплатформенное программное обеспечение для виртуализации, которое позволяет запускать несколько операционных систем (ОС) одновременно на одном физическом компьютере.
\newpage
{\centering\section{Постановка задачи}}
Установка и настройка виртуальной среды Debian 12.14, обучиться основам работы в командном интерпретаторе, создать статические и динамические библиотеки, разработка сборочных файлов CmakeLists.txt, работа с двоичными файлами и визуализация структур в псевдографике 
\newpage
{\centering\section*{Заключение}}\addcontentsline{toc}{section}{Заключение}

\newpage
\renewcommand{\bibname}{\vspace{-3.2cm}\centering\LargeСписок используемых источников}\addcontentsline{toc}{section}{Список используемых источников}
\begin{thebibliography}{8}

\bibitem{gametheory}
Петросян Л. А., Зенкевич Н. А., Семина Е. А. Теория игр : учеб. пособие для ун-тов. -- Москва : Высш. шк. : Книжный дом «Университет», 1998. - 304 с. - ISBN 5-06-001005-8.

\bibitem{virtualbox}
[Электронный ресурс] : 
\href{https://www.oracle.com/virtualization/virtualbox/}{https://www.oracle.com/virtualization/virtualbox/} 
(дата обращения: 22.05.2026).

\bibitem{habr_chess}
Habr. <<Как компьютер играет в шахматы?>> [Электронный ресурс]: 
\href{https://habr.com/ru/articles/390821/}{https://habr.com/ru/articles/390821/} 
(дата обращения: 20.05.2026).

\bibitem{habr_optimization}
Habr. <<Оптимизация перебора>> [Электронный ресурс] : 
\href{https://habr.com/ru/articles/190850/}{https://habr.com/ru/articles/190850/} 
(дата обращения: 22.05.2026).

\bibitem{devto_projects}
Adiati, A. <<Planning and Tracking Projects with GitHub's Projects Tool>> [Электронный ресурс] : 
\href{https://dev.to/adiatiayu/planning-and-tracking-projects-with-githubs-projects-tool-572l}{https://dev.to/adiatiayu/planning-and-tracking-projects-with-githubs-projects-tool-572l} 
(дата обращения: 23.05.2026).

\end{thebibliography}

\end{document}
